\uuid{jxki}
\titre{Équations complexes}
\auteur{bodin ; hanani ; mzari}
\organisation{exo7}
\licence{CC BY-NC-SA 4.0}
\datecreate{2020-01-20}
\niveau{L1}
\module{Algèbre}
\chapitre{Nombres complexes}
\sousChapitre{Racine carrée, équation du second degré}
\difficulte{2}
\theme{Équations, équations complexes}
\exo7qcmid{lille-1.nombres-complexes.023}

\contenu{
\begin{question}
\qclassification{104.02,104.03,104.04}
\qtags{motcle=équations complexes}
On considère l'équation : $(E) : \, z^5= \frac{1+i}{\sqrt 3-i}$, $z\in \Cc$.   Quelles sont les assertions vraies ?
\begin{answers}
    \bad{Si $z$ est une solution de $(E)$, $|z|=\frac{1}{\sqrt[5]{ 2}}$.}

    \good{Si $z$ est une solution de $(E)$, $|z|=\frac{1}{\sqrt[10] 2}$.}

    \bad{Si $z$ est une solution de $(E)$, $\arg z=\frac{\pi}{12} \, [2\pi]$.}

    \good{Si $z$ est une solution de $(E)$, $\arg z=\frac{\pi}{12} + \frac{2k\pi}{5} \, [2\pi], \, k \in \Zz$.}
\end{answers}
\begin{explanations}
Résoudre $ z^5= \frac{1+i}{\sqrt 3-i}= \frac{1}{\sqrt 2} e^{i\frac{5\pi}{12}}$, en utilisant l'écriture géométrique.
\end{explanations}
\end{question}
}
